\clearpage
\VerdierPhysicalPageAnchor{0062}%
\VerdierSemanticAnchor{ch1-1-6-4}
\VerdierSemanticAnchor{ch1-1-6-4-definition}
\noindent\textbf{定義 1.6.4.}\quad
\emph{$\mathbb Z^n$ の順序付き基底 $e_1,\ldots,e_n$ に関する Koszul
余サイクル}とは、基底 $e_1,\ldots,e_n$ に関する Koszul 類に属し、次を
満たすただ一つの余サイクル $\alpha$ である：

\VerdierSemanticAnchor{ch1-1-6-4-eq1}
\begin{equation}
  \alpha(\sigma_1,\sigma_2)=1
  \quad
  \forall\sigma_1\in\mathbb Z^n(\leq i),\quad
  \forall\sigma_2\in\mathbb Z^n(\geq i),\quad
  \forall i\in[1,n].
  \tag{1.6.4.1}
\end{equation}

\VerdierSemanticAnchor{ch1-1-6-4-coordinate-intro}
$\sigma_1=\sum_{i=1}^{n}p_i e_i$ および
$\sigma_2=\sum_{i=1}^{n}q_i e_i$ を $\mathbb Z^n$ の二元とすると、

\VerdierSemanticAnchor{ch1-1-6-4-eq2}
\begin{equation}
  \alpha(\sigma_1,\sigma_2)
  =(-1)^{\left(\sum_{i<j}q_i p_j\right)}.
  \tag{1.6.4.2}
\end{equation}

\VerdierSemanticAnchor{ch1-1-6-4-justification}
この定義には説明が必要である。(1.6.4.1) を満たす Koszul 類の余サイクルの
存在は明らかである。実際、直ちに計算すれば、余サイクル (1.6.4.2) は
(1.6.4.1) を満たし、Koszul 類に属することが分かる (1.6.2)。一意性を
確かめるには、$\alpha$ とコホモロガスな任意の余サイクルが

\VerdierSemanticAnchor{ch1-1-6-4-cohomologous-form}
\begin{equation*}
  \alpha_1(\sigma_1,\sigma_2)
  =\alpha(\sigma_1,\sigma_2)
   \beta(\sigma_2)\beta(\sigma_1+\sigma_2)^{-1}\beta(\sigma_1)
\end{equation*}
の形であり、この余サイクルが (1.6.4.1) を満たすことと

\VerdierSemanticAnchor{ch1-1-6-4-condition}
\begin{equation*}
  \beta(\sigma_1+\sigma_2)=\beta(\sigma_1)\beta(\sigma_2)
  \quad
  \forall i,\quad
  \forall\sigma_1\in\mathbb Z^n(\leq i),\quad
  \forall\sigma_2\in\mathbb Z^n(\geq i)
\end{equation*}
が同値であることに注意すればよい。

\VerdierSemanticAnchor{ch1-1-6-4-uniqueness}
ここから $\beta\colon\mathbb Z^n\to\mathbb Z/2\mathbb Z$ が群準同型である
ことが直ちに従い、したがって $\alpha_1=\alpha$ である。

\VerdierSemanticAnchor{ch1-1-6-5}
\VerdierSemanticAnchor{ch1-1-6-5-order-dependence}
\noindent\textbf{1.6.5.}\quad
Koszul 余サイクルの定義は、$\mathbb Z^n$ の基底 $e_1,\ldots,e_n$ に
入れた全順序に本質的に依存する。$[1,n]$ 上の自然順序を $\leq$、任意の
全順序を $\leq_s$ とする。$s$ を次を満たす $[1,n]$ のただ一つの置換とする：

\VerdierSemanticAnchor{ch1-1-6-5-eq1}
\begin{equation}
  i\leq_s j\quad\Longleftrightarrow\quad s(i)\leq s(j).
  \tag{1.6.5.1}
\end{equation}

\VerdierSemanticAnchor{ch1-1-6-5-cocycles}
$\alpha$ と $\alpha_s$ を、それぞれ順序 $\leq$ と $\leq_s$ に関する
Koszul 余サイクルとする。このとき

\VerdierSemanticAnchor{ch1-1-6-5-eq2}
\VerdierSemanticAnchor{ch1-1-6-5-eq2-natural}
\VerdierSemanticAnchor{ch1-1-6-5-eq2-coordinates}
\VerdierSemanticAnchor{ch1-1-6-5-eq2-permuted}
\begin{equation}
  \begin{aligned}
    \alpha(\sigma_1,\sigma_2)
      &=(-1)^{\left(\sum_{i<j}q_i p_j\right)},\\
    \sigma_1&=\sum_i p_i e_i,
      \qquad \sigma_2=\sum_i q_i e_i,\\
    \alpha_s(\sigma_1,\sigma_2)
      &=(-1)^{\left(\sum_{s(i)<s(j)}q_i p_j\right)},
  \end{aligned}
  \tag{1.6.5.2}
\end{equation}
\VerdierSemanticAnchor{ch1-1-6-5-open}
